\makeatletter
\newlabel{ej:baydin}{{1.5}{2}{El ejemplo de Baydin \emph {et al.}}{ejemplo.5}{}}
\newlabel{ej:baydin@cref}{{[ejemplo][5][]1.5}{[1][2][]2}}
\newlabel{fig:grafo-baydin}{{1.1}{2}{Grafo computacional de la traza del \cref {ej:baydin}. Cada nodo interior lleva la operación elemental que lo calcula}{figure.1.1}{}}
\newlabel{fig:grafo-baydin@cref}{{[figure][1][1]1.1}{[1][2][]2}}
\newlabel{ej:logistica}{{1.9}{3}{Crecimiento de las expresiones}{ejemplo.9}{}}
\newlabel{ej:logistica@cref}{{[ejemplo][9][]1.9}{[1][3][]3}}
\newlabel{eje:1-traza}{{1.1}{4}{}{ejercicio.1.1}{}}
\newlabel{eje:1-traza@cref}{{[ejercicio][1][1]1.1}{[1][4][]4}}
\newlabel{eje:1-precision}{{1.2}{4}{}{ejercicio.1.2}{}}
\newlabel{eje:1-precision@cref}{{[ejercicio][2][1]1.2}{[1][4][]4}}
\newlabel{eje:1-d2l}{{1.3}{4}{}{ejercicio.1.3}{}}
\newlabel{eje:1-d2l@cref}{{[ejercicio][3][1]1.3}{[1][4][]4}}
\newlabel{eje:1-complejo}{{1.4}{4}{}{ejercicio.1.4}{}}
\newlabel{eje:1-complejo@cref}{{[ejercicio][4][1]1.4}{[1][4][]4}}
\newlabel{eje:1-logistica}{{1.5}{4}{}{ejercicio.1.5}{}}
\newlabel{eje:1-logistica@cref}{{[ejercicio][5][1]1.5}{[1][4][]4}}
\newlabel{eje:1-grado}{{1.6}{5}{}{ejercicio.1.6}{}}
\newlabel{eje:1-grado@cref}{{[ejercicio][6][1]1.6}{[1][4][]5}}
\newlabel{ej:tangente}{{2.3}{6}{Traza tangente}{ejemplo.3}{}}
\newlabel{ej:tangente@cref}{{[ejemplo][3][]2.3}{[1][6][]6}}
\newlabel{ej:dual-cubo}{{2.10}{8}{}{ejemplo.10}{}}
\newlabel{ej:dual-cubo@cref}{{[ejemplo][10][]2.10}{[1][8][]8}}
\newlabel{eje:2-cuadrado}{{2.1}{9}{}{ejercicio.2.1}{}}
\newlabel{eje:2-cuadrado@cref}{{[ejercicio][1][2]2.1}{[1][9][]9}}
\newlabel{eje:2-cuartica}{{2.2}{9}{}{ejercicio.2.2}{}}
\newlabel{eje:2-cuartica@cref}{{[ejercicio][2][2]2.2}{[1][9][]9}}
\newlabel{eje:2-cociente}{{2.3}{9}{}{ejercicio.2.3}{}}
\newlabel{eje:2-cociente@cref}{{[ejercicio][3][2]2.3}{[1][9][]9}}
\newlabel{eje:2-tanh}{{2.4}{9}{}{ejercicio.2.4}{}}
\newlabel{eje:2-tanh@cref}{{[ejercicio][4][2]2.4}{[1][9][]9}}
\newlabel{eje:2-multidual}{{2.5}{9}{}{ejercicio.2.5}{}}
\newlabel{eje:2-multidual@cref}{{[ejercicio][5][2]2.5}{[1][9][]9}}
\newlabel{eje:2-polinomios}{{2.6}{9}{}{ejercicio.2.6}{}}
\newlabel{eje:2-polinomios@cref}{{[ejercicio][6][2]2.6}{[1][9][]9}}
\newlabel{fig:algoritmo}{{3.1}{11}{El modo inverso de la diferenciación automática~\cite [§3.2]{baydin}. En redes neuronales se llama \emph {retropropagación}~\cite [§6.5.6]{gbc}}{figure.3.1}{}}
\newlabel{fig:algoritmo@cref}{{[figure][1][3]3.1}{[1][11][]11}}
\newlabel{ej:adjunta}{{3.4}{11}{Traza adjunta del ejemplo de Baydin \emph {et al.}}{ejemplo.4}{}}
\newlabel{ej:adjunta@cref}{{[ejemplo][4][]3.4}{[1][11][]11}}
\newlabel{ej:abcf}{{3.5}{12}{Una expresión con cuatro entradas}{ejemplo.5}{}}
\newlabel{ej:abcf@cref}{{[ejemplo][5][]3.5}{[1][12][]12}}
\newlabel{fig:abcf}{{3.2}{12}{Grafo del \cref {ej:abcf}. Dentro de cada nodo, el valor de la pasada hacia delante; junto a él, en color, su adjunto, obtenido en la pasada hacia atrás}{figure.3.2}{}}
\newlabel{fig:abcf@cref}{{[figure][2][3]3.2}{[1][12][]12}}
\newlabel{ej:neurona}{{3.6}{12}{Una neurona}{ejemplo.6}{}}
\newlabel{ej:neurona@cref}{{[ejemplo][6][]3.6}{[1][12][]12}}
\newlabel{ej:orden}{{3.9}{13}{Qué ocurre si el orden es incorrecto}{ejemplo.9}{}}
\newlabel{ej:orden@cref}{{[ejemplo][9][]3.9}{[1][13][]13}}
\newlabel{eje:3-acumular}{{3.1}{14}{}{ejercicio.3.1}{}}
\newlabel{eje:3-acumular@cref}{{[ejercicio][1][3]3.1}{[1][14][]14}}
\newlabel{eje:3-producto}{{3.2}{14}{}{ejercicio.3.2}{}}
\newlabel{eje:3-producto@cref}{{[ejercicio][2][3]3.2}{[1][14][]14}}
\newlabel{eje:3-neurona}{{3.3}{14}{}{ejercicio.3.3}{}}
\newlabel{eje:3-neurona@cref}{{[ejercicio][3][3]3.3}{[1][14][]14}}
\newlabel{eje:3-topologico}{{3.4}{14}{}{ejercicio.3.4}{}}
\newlabel{eje:3-topologico@cref}{{[ejercicio][4][3]3.4}{[1][14][]14}}
\newlabel{eje:3-d2l}{{3.5}{14}{}{ejercicio.3.5}{}}
\newlabel{eje:3-d2l@cref}{{[ejercicio][5][3]3.5}{[1][14][]14}}
\newlabel{eje:3-sesgo}{{3.6}{14}{}{ejercicio.3.6}{}}
\newlabel{eje:3-sesgo@cref}{{[ejercicio][6][3]3.6}{[1][14][]14}}
\newlabel{eje:3-arbol}{{3.7}{14}{}{ejercicio.3.7}{}}
\newlabel{eje:3-arbol@cref}{{[ejercicio][7][3]3.7}{[1][14][]14}}
\newlabel{ej:vjp-componente}{{4.5}{17}{Operaciones componente a componente}{ejemplo.5}{}}
\newlabel{ej:vjp-componente@cref}{{[ejemplo][5][]4.5}{[1][17][]17}}
\newlabel{ej:vjp-lineal}{{4.7}{17}{}{ejemplo.7}{}}
\newlabel{ej:vjp-lineal@cref}{{[ejemplo][7][]4.7}{[1][17][]17}}
\newlabel{eje:4-modo}{{4.1}{18}{}{ejercicio.4.1}{}}
\newlabel{eje:4-modo@cref}{{[ejercicio][1][4]4.1}{[1][18][]18}}
\newlabel{eje:4-d2l}{{4.2}{18}{}{ejercicio.4.2}{}}
\newlabel{eje:4-d2l@cref}{{[ejercicio][2][4]4.2}{[1][18][]18}}
\newlabel{eje:4-cadena}{{4.3}{18}{}{ejercicio.4.3}{}}
\newlabel{eje:4-cadena@cref}{{[ejercicio][3][4]4.3}{[1][18][]18}}
\newlabel{eje:4-lineal}{{4.4}{18}{}{ejercicio.4.4}{}}
\newlabel{eje:4-lineal@cref}{{[ejercicio][4][4]4.4}{[1][18][]18}}
\newlabel{eje:4-cs336}{{4.5}{19}{}{ejercicio.4.5}{}}
\newlabel{eje:4-cs336@cref}{{[ejercicio][5][4]4.5}{[1][18][]19}}
\newlabel{eje:4-softmax}{{4.6}{19}{}{ejercicio.4.6}{}}
\newlabel{eje:4-softmax@cref}{{[ejercicio][6][4]4.6}{[1][19][]19}}
\newlabel{ej:hvp}{{5.5}{21}{}{ejemplo.5}{}}
\newlabel{ej:hvp@cref}{{[ejemplo][5][]5.5}{[1][21][]21}}
\newlabel{ej:recurrencia}{{5.7}{22}{Una recurrencia en el tiempo}{ejemplo.7}{}}
\newlabel{ej:recurrencia@cref}{{[ejemplo][7][]5.7}{[1][22][]22}}
\newlabel{eje:5-d2l}{{5.1}{23}{}{ejercicio.5.1}{}}
\newlabel{eje:5-d2l@cref}{{[ejercicio][1][5]5.1}{[1][22][]23}}
\newlabel{eje:5-hvp}{{5.2}{23}{}{ejercicio.5.2}{}}
\newlabel{eje:5-hvp@cref}{{[ejercicio][2][5]5.2}{[1][23][]23}}
\newlabel{eje:5-adjunto}{{5.3}{23}{}{ejercicio.5.3}{}}
\newlabel{eje:5-adjunto@cref}{{[ejercicio][3][5]5.3}{[1][23][]23}}
\newlabel{eje:5-checkpoint}{{5.4}{23}{}{ejercicio.5.4}{}}
\newlabel{eje:5-checkpoint@cref}{{[ejercicio][4][5]5.4}{[1][23][]23}}
\newlabel{eje:5-cs336}{{5.5}{23}{}{ejercicio.5.5}{}}
\newlabel{eje:5-cs336@cref}{{[ejercicio][5][5]5.5}{[1][23][]23}}
\newlabel{fig:valor}{{6.1}{25}{Núcleo de un autograd escalar: el nodo, la regla local del producto y el método \texttt {backward}, que construye un orden topológico (\cref {prop:dfs}) y lo recorre al revés}{figure.6.1}{}}
\newlabel{fig:valor@cref}{{[figure][1][6]6.1}{[1][24][]25}}
\newlabel{ej:ce-micro}{{6.3}{25}{Entropía cruzada con piezas elementales}{ejemplo.3}{}}
\newlabel{ej:ce-micro@cref}{{[ejemplo][3][]6.3}{[1][25][]25}}
\newlabel{fig:bucle}{{6.2}{26}{Un paso de entrenamiento en PyTorch: pasada hacia delante, pérdida, puesta a cero de los gradientes, modo inverso y actualización}{figure.6.2}{}}
\newlabel{fig:bucle@cref}{{[figure][2][6]6.2}{[1][26][]26}}
\newlabel{ej:mlp-caracteres}{{6.5}{26}{El modelo de caracteres del módulo~9}{ejemplo.5}{}}
\newlabel{ej:mlp-caracteres@cref}{{[ejemplo][5][]6.5}{[1][26][]26}}
\newlabel{eje:6-zero}{{6.1}{28}{}{ejercicio.6.1}{}}
\newlabel{eje:6-zero@cref}{{[ejercicio][1][6]6.1}{[1][27][]28}}
\newlabel{eje:6-dos-veces}{{6.2}{28}{}{ejercicio.6.2}{}}
\newlabel{eje:6-dos-veces@cref}{{[ejercicio][2][6]6.2}{[1][28][]28}}
\newlabel{eje:6-cmu}{{6.3}{28}{}{ejercicio.6.3}{}}
\newlabel{eje:6-cmu@cref}{{[ejercicio][3][6]6.3}{[1][28][]28}}
\newlabel{eje:6-parametros}{{6.4}{28}{}{ejercicio.6.4}{}}
\newlabel{eje:6-parametros@cref}{{[ejercicio][4][6]6.4}{[1][28][]28}}
\newlabel{eje:6-paso}{{6.5}{28}{}{ejercicio.6.5}{}}
\newlabel{eje:6-paso@cref}{{[ejercicio][5][6]6.5}{[1][28][]28}}
\newlabel{eje:6-microlotes}{{6.6}{28}{}{ejercicio.6.6}{}}
\newlabel{eje:6-microlotes@cref}{{[ejercicio][6][6]6.6}{[1][28][]28}}
\newlabel{ap:soluciones}{{A}{29}{Soluciones de los ejercicios}{appendix.A}{}}
\newlabel{ap:soluciones@cref}{{[appendix][1][2147483647]A}{[1][29][]29}}
\newlabel{sol:eje:1-traza}{{A.1}{29}{Tema 1}{section*.3}{}}
\newlabel{sol:eje:1-traza@cref}{{[subappendix][1][2147483647,1]A.1}{[1][29][]29}}
\newlabel{sol:eje:1-precision}{{A.1}{29}{Tema 1}{section*.4}{}}
\newlabel{sol:eje:1-precision@cref}{{[subappendix][1][2147483647,1]A.1}{[1][29][]29}}
\newlabel{sol:eje:1-d2l}{{A.1}{30}{Tema 1}{section*.5}{}}
\newlabel{sol:eje:1-d2l@cref}{{[subappendix][1][2147483647,1]A.1}{[1][30][]30}}
\newlabel{sol:eje:1-complejo}{{A.1}{30}{Tema 1}{section*.6}{}}
\newlabel{sol:eje:1-complejo@cref}{{[subappendix][1][2147483647,1]A.1}{[1][30][]30}}
\newlabel{sol:eje:1-logistica}{{A.1}{31}{Tema 1}{section*.7}{}}
\newlabel{sol:eje:1-logistica@cref}{{[subappendix][1][2147483647,1]A.1}{[1][30][]31}}
\newlabel{sol:eje:1-grado}{{A.1}{31}{Tema 1}{section*.8}{}}
\newlabel{sol:eje:1-grado@cref}{{[subappendix][1][2147483647,1]A.1}{[1][31][]31}}
\newlabel{sol:eje:2-cuadrado}{{A.2}{31}{Tema 2}{section*.9}{}}
\newlabel{sol:eje:2-cuadrado@cref}{{[subappendix][2][2147483647,1]A.2}{[1][31][]31}}
\newlabel{sol:eje:2-cuartica}{{A.2}{32}{Tema 2}{section*.10}{}}
\newlabel{sol:eje:2-cuartica@cref}{{[subappendix][2][2147483647,1]A.2}{[1][32][]32}}
\newlabel{sol:eje:2-cociente}{{A.2}{32}{Tema 2}{section*.11}{}}
\newlabel{sol:eje:2-cociente@cref}{{[subappendix][2][2147483647,1]A.2}{[1][32][]32}}
\newlabel{sol:eje:2-tanh}{{A.2}{32}{Tema 2}{section*.12}{}}
\newlabel{sol:eje:2-tanh@cref}{{[subappendix][2][2147483647,1]A.2}{[1][32][]32}}
\newlabel{sol:eje:2-multidual}{{A.2}{33}{Tema 2}{section*.13}{}}
\newlabel{sol:eje:2-multidual@cref}{{[subappendix][2][2147483647,1]A.2}{[1][33][]33}}
\newlabel{sol:eje:2-polinomios}{{A.2}{33}{Tema 2}{section*.14}{}}
\newlabel{sol:eje:2-polinomios@cref}{{[subappendix][2][2147483647,1]A.2}{[1][33][]33}}
\newlabel{sol:eje:3-acumular}{{A.3}{34}{Tema 3}{section*.15}{}}
\newlabel{sol:eje:3-acumular@cref}{{[subappendix][3][2147483647,1]A.3}{[1][33][]34}}
\newlabel{sol:eje:3-producto}{{A.3}{34}{Tema 3}{section*.16}{}}
\newlabel{sol:eje:3-producto@cref}{{[subappendix][3][2147483647,1]A.3}{[1][34][]34}}
\newlabel{sol:eje:3-neurona}{{A.3}{34}{Tema 3}{section*.17}{}}
\newlabel{sol:eje:3-neurona@cref}{{[subappendix][3][2147483647,1]A.3}{[1][34][]34}}
\newlabel{sol:eje:3-topologico}{{A.3}{35}{Tema 3}{section*.18}{}}
\newlabel{sol:eje:3-topologico@cref}{{[subappendix][3][2147483647,1]A.3}{[1][35][]35}}
\newlabel{sol:eje:3-d2l}{{A.3}{35}{Tema 3}{section*.19}{}}
\newlabel{sol:eje:3-d2l@cref}{{[subappendix][3][2147483647,1]A.3}{[1][35][]35}}
\newlabel{sol:eje:3-sesgo}{{A.3}{35}{Tema 3}{section*.20}{}}
\newlabel{sol:eje:3-sesgo@cref}{{[subappendix][3][2147483647,1]A.3}{[1][35][]35}}
\newlabel{sol:eje:3-arbol}{{A.3}{36}{Tema 3}{section*.21}{}}
\newlabel{sol:eje:3-arbol@cref}{{[subappendix][3][2147483647,1]A.3}{[1][36][]36}}
\newlabel{sol:eje:4-modo}{{A.4}{36}{Tema 4}{section*.22}{}}
\newlabel{sol:eje:4-modo@cref}{{[subappendix][4][2147483647,1]A.4}{[1][36][]36}}
\newlabel{sol:eje:4-d2l}{{A.4}{37}{Tema 4}{section*.23}{}}
\newlabel{sol:eje:4-d2l@cref}{{[subappendix][4][2147483647,1]A.4}{[1][37][]37}}
\newlabel{sol:eje:4-cadena}{{A.4}{37}{Tema 4}{section*.24}{}}
\newlabel{sol:eje:4-cadena@cref}{{[subappendix][4][2147483647,1]A.4}{[1][37][]37}}
\newlabel{sol:eje:4-lineal}{{A.4}{38}{Tema 4}{section*.25}{}}
\newlabel{sol:eje:4-lineal@cref}{{[subappendix][4][2147483647,1]A.4}{[1][38][]38}}
\newlabel{sol:eje:4-cs336}{{A.4}{38}{Tema 4}{section*.26}{}}
\newlabel{sol:eje:4-cs336@cref}{{[subappendix][4][2147483647,1]A.4}{[1][38][]38}}
\newlabel{sol:eje:4-softmax}{{A.4}{38}{Tema 4}{section*.27}{}}
\newlabel{sol:eje:4-softmax@cref}{{[subappendix][4][2147483647,1]A.4}{[1][38][]38}}
\newlabel{sol:eje:5-d2l}{{A.5}{39}{Tema 5}{section*.28}{}}
\newlabel{sol:eje:5-d2l@cref}{{[subappendix][5][2147483647,1]A.5}{[1][39][]39}}
\newlabel{sol:eje:5-hvp}{{A.5}{39}{Tema 5}{section*.29}{}}
\newlabel{sol:eje:5-hvp@cref}{{[subappendix][5][2147483647,1]A.5}{[1][39][]39}}
\newlabel{sol:eje:5-adjunto}{{A.5}{40}{Tema 5}{section*.30}{}}
\newlabel{sol:eje:5-adjunto@cref}{{[subappendix][5][2147483647,1]A.5}{[1][40][]40}}
\newlabel{sol:eje:5-checkpoint}{{A.5}{40}{Tema 5}{section*.31}{}}
\newlabel{sol:eje:5-checkpoint@cref}{{[subappendix][5][2147483647,1]A.5}{[1][40][]40}}
\newlabel{sol:eje:5-cs336}{{A.5}{41}{Tema 5}{section*.32}{}}
\newlabel{sol:eje:5-cs336@cref}{{[subappendix][5][2147483647,1]A.5}{[1][40][]41}}
\newlabel{sol:eje:6-zero}{{A.6}{41}{Tema 6}{section*.33}{}}
\newlabel{sol:eje:6-zero@cref}{{[subappendix][6][2147483647,1]A.6}{[1][41][]41}}
\newlabel{sol:eje:6-dos-veces}{{A.6}{42}{Tema 6}{section*.34}{}}
\newlabel{sol:eje:6-dos-veces@cref}{{[subappendix][6][2147483647,1]A.6}{[1][41][]42}}
\newlabel{sol:eje:6-cmu}{{A.6}{42}{Tema 6}{section*.35}{}}
\newlabel{sol:eje:6-cmu@cref}{{[subappendix][6][2147483647,1]A.6}{[1][42][]42}}
\newlabel{sol:eje:6-parametros}{{A.6}{42}{Tema 6}{section*.36}{}}
\newlabel{sol:eje:6-parametros@cref}{{[subappendix][6][2147483647,1]A.6}{[1][42][]42}}
\newlabel{sol:eje:6-paso}{{A.6}{43}{Tema 6}{section*.37}{}}
\newlabel{sol:eje:6-paso@cref}{{[subappendix][6][2147483647,1]A.6}{[1][43][]43}}
\newlabel{sol:eje:6-microlotes}{{A.6}{43}{Tema 6}{section*.38}{}}
\newlabel{sol:eje:6-microlotes@cref}{{[subappendix][6][2147483647,1]A.6}{[1][43][]43}}
\makeatother
